\documentclass{beamer}

\begin{document}

\begin{frame}
\frametitle{Problem 1: Elasticity of Gig Economy Employment to COVID-19}
\textbf{Problem Statement:}
You are tasked with calculating the \textbf{elasticity of gig economy employment} with respect to COVID-19 cases. The elasticity will help measure the percentage change in gig economy workers for each 1\% change in COVID-19 cases.

\textbf{Tasks:}
\begin{enumerate}
    \item \textbf{Elasticity Calculation:}
    \item \textbf{Interpretation:}
    \item \textbf{Impact of Increasing COVID-19 Cases:}
\end{enumerate}
\end{frame}

\begin{frame}
\frametitle{Solution to Problem 1: Elasticity of Gig Economy Employment to COVID-19}
\textbf{1. Elasticity Calculation:}

The formula for elasticity is:

\[
\text{Elasticity} = \frac{\% \Delta \text{Gig Workers}}{\% \Delta \text{COVID-19 Cases}}
\]

Where:
\[
\% \Delta \text{Gig Workers} = \frac{\text{Gig Workers}_t - \text{Gig Workers}_{t-1}}{\text{Gig Workers}_{t-1}} \times 100
\]
\[
\% \Delta \text{COVID-19 Cases} = \frac{\text{COVID-19 Cases}_t - \text{COVID-19 Cases}_{t-1}}{\text{COVID-19 Cases}_{t-1}} \times 100
\]

Elasticity measures how sensitive gig economy employment is to changes in COVID-19 cases. A high elasticity indicates that a small change in COVID-19 cases leads to a large change in gig economy employment.

\end{frame}

\begin{frame}
\frametitle{Solution to Problem 1: Elasticity of Gig Economy Employment to COVID-19 (cont.)}
\textbf{2. Interpretation:}

If the elasticity is greater than 1:
\[
\text{Elasticity} > 1 \quad \Rightarrow \quad \text{Elastic Response}
\]
This implies that gig economy employment is highly responsive to changes in COVID-19 cases. A 1\% increase in COVID-19 cases results in a greater than 1\% decrease in gig economy workers.

If the elasticity is less than 1:
\[
\text{Elasticity} < 1 \quad \Rightarrow \quad \text{Inelastic Response}
\]
This implies that gig economy employment is less responsive to changes in COVID-19 cases. A 1\% increase in COVID-19 cases leads to a less than 1\% change in gig economy workers.

\end{frame}

\begin{frame}
\frametitle{Solution to Problem 1: Elasticity of Gig Economy Employment to COVID-19 (cont.)}
\textbf{3. Impact of Increasing COVID-19 Cases:}

If COVID-19 cases increase, we expect the gig economy to react differently based on the sector and the elasticity of employment. For sectors such as transportation (e.g., Uber, Lyft), which depend on physical mobility, the response may be highly elastic, leading to a significant reduction in gig workers. Conversely, sectors like freelance work (e.g., Upwork, Fiverr) may see a less elastic response, with fewer changes in employment levels.

As COVID-19 cases rise, the gig economy may face greater challenges, but the degree of sensitivity depends on the type of work and the reliance on in-person services.

\end{frame}

\begin{frame}
\frametitle{Problem 2: COVID-19’s Differential Impact on Gig Economy Sectors}
\textbf{Problem Statement:}
The gig economy consists of various sectors such as transportation, food delivery, and freelance work. The pandemic may have affected each of these sectors differently due to differences in demand, restrictions, and labor conditions.

\textbf{Tasks:}
\begin{enumerate}
    \item \textbf{Sectoral Differences:}
    \item \textbf{Policy Implications:}
    \item \textbf{Long-Term Sectoral Shifts:}
\end{enumerate}
\end{frame}

\begin{frame}
\frametitle{Solution to Problem 2: COVID-19’s Differential Impact on Gig Economy Sectors}
\textbf{1. Sectoral Differences:}

The impact of COVID-19 on gig economy sectors varies due to factors like:
\begin{itemize}
    \item \textbf{Transportation (e.g., Uber, Lyft):} Highly dependent on physical movement and face-to-face interaction. With restrictions on mobility, lockdowns, and social distancing, the demand for transportation services decreased, leading to a significant reduction in gig economy workers.
    \item \textbf{Food Delivery (e.g., SkipTheDishes, DoorDash):} Demand for food delivery services increased as restaurants shifted to takeout and delivery models. Therefore, gig workers in this sector experienced a rise in employment as consumer demand for delivery services grew during lockdowns.
    \item \textbf{Freelance Work (e.g., Upwork, Fiverr):} With more businesses adopting remote work models, freelance opportunities in digital services surged. As a result, this sector may have seen a rise in gig workers as companies and individuals sought out freelancers for virtual work.
\end{itemize}
\end{frame}

\begin{frame}
\frametitle{Solution to Problem 2: COVID-19’s Differential Impact on Gig Economy Sectors (cont.)}
\textbf{2. Policy Implications:}

Governments and digital platforms can introduce several policies to mitigate the negative effects of COVID-19 on specific gig economy sectors:
\begin{itemize}
    \item \textbf{Transportation:} Governments could provide financial support or subsidies to gig workers in the transportation sector. Digital platforms could introduce measures like contactless payments and health guidelines to protect workers and consumers.
    \item \textbf{Food Delivery:} Platforms could enhance safety measures, provide financial support for delivery workers, and incentivize consumers to use delivery services to support gig workers. Temporary income support could be provided during periods of lockdown.
    \item \textbf{Freelance Work:} Governments could offer financial relief to freelancers through grants or subsidies. Digital platforms could ensure job security by providing more projects to freelancers in growing sectors such as virtual assistance and IT support.
\end{itemize}
\end{frame}

\begin{frame}
\frametitle{Solution to Problem 2: COVID-19’s Differential Impact on Gig Economy Sectors (cont.)}
\textbf{3. Long-Term Sectoral Shifts:}

The pandemic may cause long-term shifts in the types of gig economy work:
\begin{itemize}
    \item \textbf{Shift towards Digital Freelance Opportunities:} As more businesses embrace remote work, there could be an increased demand for digital freelancers (e.g., graphic designers, software developers, writers) and a decline in demand for in-person services in sectors like transportation.
    \item \textbf{Changes in Mobility-Dependent Sectors:} As mobility restrictions ease, gig workers in transportation sectors may see a recovery in demand. However, some consumers may have developed preferences for other modes of transportation, or may opt for permanent shifts to remote working, affecting long-term demand.
    \item \textbf{Greater Use of Automation:} Automation technologies, such as autonomous vehicles or AI-powered customer service, could reduce the demand for gig workers in certain sectors, such as transportation and customer service, in the longer run.
\end{itemize}
\end{frame}

\end{document}
